\chapter{2015年广东卷（理）}

\section{单选题}

\begin{problem}
若集合 \(M = \{x \mid (x + 4)(x + 1) = 0\}\)，\(N = \{x \mid (x - 4)(x - 1) = 0\}\)，则 \(M \cap N =\)（\quad）

\choices
    {\(\{1, 4\}\)}
    {\( \{-1,-4\} \)}
    {0}
    {\(\varnothing\)}
\end{problem}

\begin{problem}
若复数 \(z = \i(3 - 2\i)\) (i是虚数单位)，则 \(z =\)（\quad）

\choices
    {\(2 - 3\i\)}
    {\(2 + 3\i\)}
    {\(3 + 2\i\)}
    {\(3 - 2\i\)}
\end{problem}

\begin{problem}
下列函数中，既不是奇函数，也不是偶函数的是（\quad）

\choices
    {\( y = \sqrt{1 + x^2} \)}
    {\( y = x + \frac{1}{x} \)}
    {\( y = 2^x + \frac{1}{2^x} \)}
    {\(y = x + \e^{x}\)}
\end{problem}

\begin{problem}
袋中共有 15 个除了颜色外完全相同 的球，其中有 10 个白球，5 个红球. 从 袋中任取2个球，所取的2个球中恰有1个白球，1个红球的概率为（\quad）

\choices
    {\( \frac{3}{91} \)}
    {\( \frac{10}{91} \)}
    {\( \frac{11}{91} \)}
    {1}
\end{problem}

\begin{problem}
平行于直线 \(2x + y + 1 = 0\) 且与圆 \(x^{2} + y^{2} = 5\) 相切的直线的方程是（\quad）

\choices
    {\(2x + y + 5 = 0\) 或 \(2x + y - 5 = 0\)}
    {\(2x + y + \sqrt{5} = 0\) 或 \(2x + y - \sqrt{5} = 0\)}
    {\(2x - y + 5 = 0\) 或 \(2x - y - 5 = 0\)}
    {\(2x - y + \sqrt{5} = 0\) 或 \(2x - y - \sqrt{5} = 0\)}
\end{problem}

\begin{problem}
若变量 \(x, y\) 满足约束条件 \(\left\{ \begin{array}{l} 4x + 5y \geqslant 8, \\ 1 \leqslant z \leqslant 3, \\ 0 \leqslant y \leqslant 2, \end{array} \right.\) 则 \(z = 3x + 2y\) 的最小值为（\quad）

\choices
    {4}
    {\(\frac{23}{5}\)}
    {6}
    {\(\frac{31}{5}\)}
\end{problem}

\begin{problem}
已知双曲线 \( C: \frac{x^{2}}{a^{2}} - \frac{y^{2}}{b^{2}} = 1 \) 的离心率 \( e = \frac{5}{4} \) ，且其右焦点为 \( F_{2}(5,0) \) ，则双曲线 C 的方程为（\quad）

\choices
    {\(\frac{x^2}{4} -\frac{y^2}{3} = 1\)}
    {\(\frac{x^2}{9} -\frac{y^2}{16} = 1\)}
    {\(\frac{x^2}{16} -\frac{y^2}{9} = 1\)}
    {\(\frac{x^2}{3} -\frac{y^2}{4} = 1\)}
\end{problem}

\begin{problem}
若空间中 \(n\) 个不同的点两两距离都相等，则正整数 \(n\) 的取值（\quad）

\choices
    {至多等于 3}
    {至多等于 4}
    {等于 5}
    {大于 5}
\end{problem}

\section{填空题}

\begin{problem}
在 \((\sqrt{a} - 1)^4\) 的展开式中，\(x\) 的系数为 \fillinblank{} \fillinblank{}.
\end{problem}

\begin{problem}
在等差数列 \(\{n_n\}\) 中，若 \(a_3 + a_4 + a_5 + a_6 + a_7 = 25\)，则 \(a_2 + a_8 = \)  \fillinblank{}.
\end{problem}

\begin{problem}
设 \(\triangle ABC\) 的内角 \(A, B, C\) 的对边分别为 \(a, b, c\). 若 \(a = \sqrt{3}, \sin B = \frac{1}{2}, C = \frac{5}{6}\)，则 \(b = \)  \fillinblank{}.
\end{problem}

\begin{problem}
某高三毕业班有 40 人，同学之间两两彼此给对方仅写一条毕业留言，那么全班共写了条毕业留言. (用数字作答)
\end{problem}

\begin{problem}
已知随机变量 X 服从二项分布 \( B(n,p) \) ，若 \( E(X)=30, D(X)=20 \) ，则 p=\fillinblank{} \fillinblank{}.
\end{problem}

\begin{problem}
已知直线 l 的极坐标方程为 \( 2\rho\sin\left(\theta-\frac{\pi}{4}\right)=\sqrt{2} \) ，点 A 的极坐标为 \( A\left(2\sqrt{2},\frac{7\pi}{4}\right) \) ，则点 A 到直线 l 的距离为 \fillinblank{} \fillinblank{}.
\end{problem}

\begin{problem}
如图，已知 \(AB\) 是圆 \(O\) 的直径，\(AB = 4\)，\(EC\) 是圆 \(O\) 的切线，切点为 \(C\)，\(BC = 1\). 过圆心 \(O\) 作 \(BC\) 的平行线，分别交 \(EC\) 和 \(AC\) 于点 \(D\) 和点 \(P\)，则 \(OD = \fillinblank{}\).
\end{problem}

\section{解答题}

\begin{problem}
在平面直角坐标系 \(xOy\) 中，已知向量 \(\bs{m} = \begin{pmatrix} \sqrt{2} & \sqrt{2} \\ 2 & -2 \end{pmatrix}\)，\(\bs{n} = (\sin x, \cos x), x \in \left(0, \frac{\pi}{2}\right)\). (1) 若 \(m \perp n\)，求 \(\tan x\) 的值; (2) 若 \(m\) 与 \(n\) 的夹角为 \(\frac{5}{3}\)，求 \(x\) 的值.
\end{problem}

\begin{problem}
某工厂 36 名工人的年龄数据如下表. 工人编号 年龄 工人编号 年龄 工人编号 年龄 1 40 10 36 19 27 2 44 11 31 20 43 3 40 12 38 21 41 4 41 13 39 22 37 5 33 14 43 23 31 6 40 15 45 24 42 7 45 16 39 25 37 8 42 17 38 26 34 9 43 18 36 27 42 (1) 用系统抽样法从 36 名工人中抽取容量为 9 样的，且在第一分段里用随机抽样法抽到的年龄数据为 44，列出样本的年龄数据; (2) 计算 (1) 中样本的均值 \(\overline{x}\) 和方差 \(s^2\); (3) 36 名工人中年龄在 \(\overline{x} - s\) 与 \(\overline{x} + s\) 之间有多少人？所占的百分比是多少（精确到 \(0.01\%\)）\(\perp\) 
\end{problem}

\begin{problem}
如图，三角形 \(PDC\) 所在的平面与长方形 \(ABCD\) 所在的平面垂直，\(PD = PC = 4\)，\(AB = 6\)，\(BC = 3\). 点 \(E\) 是 \(CD\) 边的中点，点 \(F\)，\(G\) 分别在线段 \(AB\)，\(BC\) 上，且 \(AF = 2FB\)，\(CG = 2GB\). (1) 证明: \(PE \perp FG\); (2) 求二面角 \(P - AD - C\) 的正切值； (3) 求直线 PA 与直线 FG 所成角的余弦值.
\end{problem}

\begin{problem}
已知过原点的动直线 \(l\) 与圆 \(C_1: x^2 + y^2 - 6x + 5 = 0\) 相交于不同的两点 \(A, B\). (1) 求圆 \(C_1\) 的圆心坐标; (2) 求线段 \(AB\) 的中点 \(M\) 的轨迹 \(C\) 的方程; (3) 是否存在实数 \(k\)，使得直线 \(L: y = k(x - 4)\) 与曲线 \(C\) 只有一个交点\(\perp\) 若存在，求出 \(k\) 的取值范围. 若不存在，说明理由.
\end{problem}

\begin{problem}
数列 \(\{a_{n}\}\) 满足 \(a_1 + 2a_2 + \dots + na_n = 4 - \frac{n + 2}{2^{n - 1}}, n \in \mathbf{N}^*\). (1) 求 \(a_{3}\) 的值; (2) 求数列 \( \{a_{n}\} \) 的前 n 项和 \( T_{n} \) ; (3) 令 \( b_{1} = a_{1}, b_{n} = \frac{T_{n - 1}}{n} + \left(1 + \frac{1}{2} + \frac{1}{3} + \dots + \frac{1}{n}\right) a_{n} (n \geqslant 2) \)，证明: 数列 \( \{b_{n}\} \) 的前 \( n \) 项和 \( S_{n} \) 满足 \( S_{n} < 2 + 2\ln n \).
\end{problem}

\begin{problem}
设 \(a > 1\)，函数 \(f(x) = (1 + x^2)\e^x - a\). (1) 求 \( f(x) \) 的单调区间: (2) 证明: \(f(x)\) 在 \((- \infty, + \infty)\) 上仅有一个零点; (3) 若曲线 \( y = f(x) \) 在点 \( P \) 处的切线与 \( x \) 轴平行，且在点 \( M(m, n) \) 处的切线与直线 \( OP \) 平行 (\( O \) 是坐标原点)，证明: \( m \leqslant \sqrt[n]{a - \frac{2}{\e}} - 1 \).
\end{problem}

